\documentclass[a4paper,12pt]{paper}

\usepackage{mycommands}

% Worked-example layout: boxed problem statement followed by a solution.
% (paper.cls already defines a list environment called example, hence \renewenvironment.)
\newcounter{example}
\renewenvironment{example}[1]{%
  \refstepcounter{example}%
  \begin{mdframed}[linewidth=1pt,innertopmargin=6pt,innerbottommargin=6pt]%
  \noindent\textbf{Example \theexample: #1.}\ }{%
  \end{mdframed}}
\newenvironment{solution}{\par\noindent\textbf{Solution.}\ }{\par\medskip}

\title{Worked examples for Lecture 13: Constitutive theory}
\author{David Al-Attar \\ Michaelmas Term, \the\year}

\begin{document}

\maketitle

\noindent These examples accompany Lecture~13. They are intended as practice in the concrete
calculations that underlie the lecture -- differentiating strain energy functions, linearising
about a stress-free equilibrium, and handling the symmetries of the elastic tensor -- and are
less demanding than the questions on the problem sets. Numerical results were checked with the
scripts in the \texttt{python} directory. Throughout we write $H_{ij}=\partial u_{i}/\partial x_{j}$
for the displacement gradient of a linearised motion and
$e_{ij}=\frac{1}{2}(H_{ij}+H_{ji})$ for the linearised strain tensor.

\begin{example}{The St Venant--Kirchhoff material}
The simplest strain energy function compatible with material frame indifference and isotropy is
that of the St Venant--Kirchhoff material,
\begin{equation}
U(\mathbf{C})=\frac{1}{2}\lambda\,(\mathrm{tr}\,\mathbf{E})^{2}+\mu\,\mathrm{tr}(\mathbf{E}^{2}),\quad
\mathbf{E}=\frac{1}{2}(\mathbf{C}-\mathbf{I}),
\label{eq:1}
\end{equation}
where $\lambda$ and $\mu$ are constants and $\mathbf{E}$ is the finite strain tensor.
(a) Compute the second Piola--Kirchhoff stress $\mathbf{S}=2\,\partial U/\partial\mathbf{C}$ and the
first Piola--Kirchhoff stress $\mathbf{T}=\mathbf{F}\mathbf{S}$, and verify that $\mathbf{F}^{-1}\mathbf{T}$
is symmetric. (b) Linearise about $\mathbf{F}=\mathbf{I}$ by writing $\mathbf{F}=\mathbf{I}+s\mathbf{H}$
and show that
$T_{ij}=s\left[\lambda\delta_{ij}\delta_{kl}+\mu(\delta_{ik}\delta_{jl}+\delta_{il}\delta_{jk})\right]H_{kl}+O(s^{2})$,
so that the elastic tensor is the isotropic one quoted in the lecture. (c) Check the result by
expanding $W(\mathbf{F})=U(\mathbf{F}^{T}\mathbf{F})$ directly to second order in $s$.
\end{example}

\begin{solution}
(a) The chain-rule argument of the lecture treats the nine components $C_{kl}$ as independent
variables, so we must regard eq.~(\ref{eq:1}) as defining $U$ on all $3\times3$ matrices, with
$\mathrm{tr}(\mathbf{E}^{2})=E_{kl}E_{lk}$. This extension is invariant under
$\mathbf{C}\rightarrow\mathbf{C}^{T}$, which is what the lecture's derivation of
$\mathbf{T}=2\mathbf{F}\,\partial U/\partial\mathbf{C}$ requires. From $\partial E_{kl}/\partial C_{ij}
=\frac{1}{2}\delta_{ki}\delta_{lj}$ we have
\begin{equation}
\frac{\partial\,\mathrm{tr}\,\mathbf{E}}{\partial C_{ij}}=\frac{1}{2}\delta_{ij},\quad
\frac{\partial}{\partial C_{ij}}\left(E_{kl}E_{lk}\right)=2E_{lk}\frac{\partial E_{kl}}{\partial C_{ij}}=E_{ji},
\label{eq:2}
\end{equation}
and hence, for symmetric $\mathbf{C}$,
\begin{equation}
S_{ij}=2\frac{\partial U}{\partial C_{ij}}=\lambda\,(\mathrm{tr}\,\mathbf{E})\,\delta_{ij}+2\mu E_{ij},\quad
\mathbf{T}=\mathbf{F}\mathbf{S}=\mathbf{F}\left[\lambda\,(\mathrm{tr}\,\mathbf{E})\,\mathbf{I}+2\mu\mathbf{E}\right].
\label{eq:3}
\end{equation}
The stress--strain relation between $\mathbf{S}$ and $\mathbf{E}$ therefore has exactly the form of
Hooke's law, but here it holds for finite deformations. Since $\mathbf{F}^{-1}\mathbf{T}=\mathbf{S}$
and $\mathbf{E}$ is symmetric, $\mathbf{F}^{-1}\mathbf{T}$ is symmetric, as material frame
indifference requires. Note that $\mathbf{T}$ itself is not symmetric: for the simple shear
$\mathbf{F}=\mathbf{I}+g\,\hat{\mathbf{e}}_{1}\hat{\mathbf{e}}_{2}^{T}$ a short calculation from
eq.~(\ref{eq:3}) gives $T_{21}=\mu g$ but $T_{12}=\mu g+g^{3}(\tfrac{1}{2}\lambda+\mu)$.

(b) With $\mathbf{F}=\mathbf{I}+s\mathbf{H}$,
\begin{equation}
\mathbf{C}=\mathbf{I}+s(\mathbf{H}+\mathbf{H}^{T})+s^{2}\mathbf{H}^{T}\mathbf{H},\quad
\mathbf{E}=s\,\mathbf{e}+\frac{1}{2}s^{2}\mathbf{H}^{T}\mathbf{H},
\label{eq:4}
\end{equation}
so the finite strain tensor reduces to the linearised strain at first order. Substituting into
eq.~(\ref{eq:3}) and discarding $O(s^{2})$ terms, including those generated by the factor
$\mathbf{F}=\mathbf{I}+s\mathbf{H}$ multiplying $\mathbf{S}$,
\begin{equation}
T_{ij}=s\left[\lambda\,e_{kk}\,\delta_{ij}+2\mu\,e_{ij}\right]+O(s^{2})
=s\left[\lambda\delta_{ij}\delta_{kl}+\mu(\delta_{ik}\delta_{jl}+\delta_{il}\delta_{jk})\right]H_{kl}+O(s^{2}),
\label{eq:5}
\end{equation}
where in the second step we used $2e_{ij}=H_{ij}+H_{ji}=(\delta_{ik}\delta_{jl}+\delta_{il}\delta_{jk})H_{kl}$.
Because $\mathbf{H}$ is an arbitrary matrix, not necessarily symmetric, the coefficient of $H_{kl}$
is uniquely determined, and comparison with Hooke's law $T_{ij}=sA_{ijkl}\,\partial u_{k}/\partial x_{l}$
identifies
\begin{equation}
A_{ijkl}=\lambda\delta_{ij}\delta_{kl}+\mu(\delta_{ik}\delta_{jl}+\delta_{il}\delta_{jk}),
\label{eq:6}
\end{equation}
which is the isotropic elastic tensor of the lecture, with $\lambda$ and $\mu$ the Lam\'{e} parameters.

(c) Since $W(\mathbf{I})=0$ and $\mathbf{T}(\mathbf{I})=\mathbf{0}$, Taylor's theorem gives
\begin{equation}
W(\mathbf{I}+s\mathbf{H})=\frac{1}{2}s^{2}A_{ijkl}H_{ij}H_{kl}+O(s^{3}),\quad
A_{ijkl}=\frac{\partial^{2}W}{\partial F_{ij}\partial F_{kl}}(\mathbf{I}).
\label{eq:7}
\end{equation}
On the other hand, from
eq.~(\ref{eq:4}), $(\mathrm{tr}\,\mathbf{E})^{2}=s^{2}(e_{kk})^{2}+O(s^{3})$ and
$\mathrm{tr}(\mathbf{E}^{2})=s^{2}e_{ij}e_{ij}+O(s^{3})$, so that
\begin{equation}
W(\mathbf{I}+s\mathbf{H})=s^{2}\left[\frac{1}{2}\lambda\,(e_{kk})^{2}+\mu\,e_{ij}e_{ij}\right]+O(s^{3}).
\label{eq:8}
\end{equation}
Using $e_{ij}e_{ij}=\frac{1}{2}(H_{ij}H_{ij}+H_{ij}H_{ji})$ one checks that the bracket equals
$\frac{1}{2}A_{ijkl}H_{ij}H_{kl}$ with $A_{ijkl}$ as in eq.~(\ref{eq:6}), as required. Writing
$e_{ij}=e'_{ij}+\frac{1}{3}e_{kk}\delta_{ij}$ with $e'_{ij}$ traceless, the quadratic form becomes
\begin{equation}
\frac{1}{2}A_{ijkl}H_{ij}H_{kl}=\frac{1}{2}\kappa\,(e_{kk})^{2}+\mu\,e'_{ij}e'_{ij},\quad
\kappa=\lambda+\frac{2}{3}\mu,
\label{eq:9}
\end{equation}
so the requirement that the stress-free configuration be a minimum of the elastic energy is
$\kappa>0$ and $\mu>0$: the bulk modulus of the lecture controls the energy of volume changes, and
the shear modulus that of shape changes.
\end{solution}

\begin{example}{An elastic fluid and the Cauchy stress}
An elastic fluid has strain energy function $W(\mathbf{F})=V(J)$ with $J=\det\mathbf{F}$.
(a) Using the cofactor expansion of the determinant, show that
$\partial J/\partial F_{ij}=J\,(\mathbf{F}^{-1})_{ji}$, i.e.\ $\partial J/\partial\mathbf{F}=J\mathbf{F}^{-T}$.
(b) Deduce that $\mathbf{T}=V'(J)\,J\,\mathbf{F}^{-T}$ and check that $\mathbf{F}^{-1}\mathbf{T}$ is symmetric.
(c) Let $\dd S$ be a surface element of the reference body with unit normal $\hat{\mathbf{n}}$, and
$\dd S_{t}$, $\hat{\mathbf{n}}_{t}$ its image under the motion. The Cauchy stress $\bm{\sigma}$ is
defined so that the force on the element is $\bm{\sigma}\hat{\mathbf{n}}_{t}\dd S_{t}=\mathbf{T}\hat{\mathbf{n}}\dd S$.
Show that $\hat{\mathbf{n}}_{t}\dd S_{t}=J\mathbf{F}^{-T}\hat{\mathbf{n}}\dd S$, and hence that
$\bm{\sigma}=J^{-1}\mathbf{T}\mathbf{F}^{T}$. Deduce that the Cauchy stress in the fluid is
$-p\mathbf{I}$ with $p=-V'(J)$. (d) Linearise about a stress-free equilibrium, $V'(1)=0$, and show
that $A_{ijkl}=\kappa\delta_{ij}\delta_{kl}$ with $\kappa=V''(1)$.
\end{example}

\begin{solution}
(a) Expanding the determinant along the $i$th row,
\begin{equation}
J=\sum_{j=1}^{3}F_{ij}\,\mathrm{cof}_{ij}(\mathbf{F})\quad\text{(no sum on $i$)},
\label{eq:10}
\end{equation}
where the cofactor $\mathrm{cof}_{ij}(\mathbf{F})$ is $(-1)^{i+j}$ times the determinant of the
$2\times2$ matrix obtained by deleting row $i$ and column $j$. Since no element of row $i$ appears in
any $\mathrm{cof}_{ij}(\mathbf{F})$, differentiating eq.~(\ref{eq:10}) with respect to
$F_{ij}$ gives $\partial J/\partial F_{ij}=\mathrm{cof}_{ij}(\mathbf{F})$. Cramer's rule states that
$\mathbf{F}^{-1}=J^{-1}\mathrm{cof}(\mathbf{F})^{T}$, so $\mathrm{cof}_{ij}(\mathbf{F})=J(\mathbf{F}^{-1})_{ji}$ and
\begin{equation}
\frac{\partial J}{\partial F_{ij}}=J\,(\mathbf{F}^{-1})_{ji}=J\,(\mathbf{F}^{-T})_{ij}.
\label{eq:11}
\end{equation}
At $\mathbf{F}=\mathbf{I}$ this reads $\partial J/\partial F_{ij}=\delta_{ij}$, i.e.\
$\det(\mathbf{I}+s\mathbf{H})=1+s\,\mathrm{tr}\,\mathbf{H}+O(s^{2})$, which we use in part (d).

(b) By the chain rule and eq.~(\ref{eq:11}),
\begin{equation}
T_{ij}=\frac{\partial W}{\partial F_{ij}}=V'(J)\frac{\partial J}{\partial F_{ij}}=V'(J)\,J\,(\mathbf{F}^{-T})_{ij},
\label{eq:12}
\end{equation}
and $\mathbf{F}^{-1}\mathbf{T}=V'(J)\,J\,\mathbf{F}^{-1}\mathbf{F}^{-T}=V'(J)\,J\,\mathbf{C}^{-1}$ is
symmetric; equivalently the second Piola--Kirchhoff stress is $\mathbf{S}=V'(J)\,J\,\mathbf{C}^{-1}$.
The first Piola--Kirchhoff stress $\mathbf{T}$ is proportional to $\mathbf{F}^{-T}$ and is
symmetric only when $\mathbf{F}$ is.

(c) A surface element at $\mathbf{x}$ spanned by the small vectors $\dd\mathbf{a}$ and $\dd\mathbf{b}$
has $\hat{\mathbf{n}}\dd S=\dd\mathbf{a}\times\dd\mathbf{b}$, and under the motion the spanning
vectors become $\mathbf{F}\dd\mathbf{a}$ and $\mathbf{F}\dd\mathbf{b}$, so
$\hat{\mathbf{n}}_{t}\dd S_{t}=(\mathbf{F}\dd\mathbf{a})\times(\mathbf{F}\dd\mathbf{b})$. For any
vectors $\mathbf{a}$, $\mathbf{b}$, $\mathbf{c}$ the scalar triple product is a determinant, so
\begin{equation}
\left[(\mathbf{F}\mathbf{a})\times(\mathbf{F}\mathbf{b})\right]\cdot(\mathbf{F}\mathbf{c})
=\det\left(\mathbf{F}\left[\mathbf{a}\;\mathbf{b}\;\mathbf{c}\right]\right)
=J\,(\mathbf{a}\times\mathbf{b})\cdot\mathbf{c},
\label{eq:13}
\end{equation}
where $[\mathbf{a}\;\mathbf{b}\;\mathbf{c}]$ is the matrix with the given columns. The left-hand
side equals $\mathbf{F}^{T}[(\mathbf{F}\mathbf{a})\times(\mathbf{F}\mathbf{b})]\cdot\mathbf{c}$, and
as $\mathbf{c}$ is arbitrary,
\begin{equation}
(\mathbf{F}\mathbf{a})\times(\mathbf{F}\mathbf{b})=J\mathbf{F}^{-T}(\mathbf{a}\times\mathbf{b}),
\quad\text{hence}\quad
\hat{\mathbf{n}}_{t}\dd S_{t}=J\mathbf{F}^{-T}\hat{\mathbf{n}}\dd S,
\label{eq:14}
\end{equation}
which is Nanson's formula; the matrix $J\mathbf{F}^{-T}$ is the cofactor matrix of part (a), which
is no coincidence, as cofactors are signed areas. Substituting into the definition of the Cauchy
stress, $\bm{\sigma}J\mathbf{F}^{-T}\hat{\mathbf{n}}=\mathbf{T}\hat{\mathbf{n}}$ for every
$\hat{\mathbf{n}}$, and therefore
\begin{equation}
\bm{\sigma}=J^{-1}\mathbf{T}\mathbf{F}^{T}=J^{-1}\mathbf{F}\left(\mathbf{F}^{-1}\mathbf{T}\right)\mathbf{F}^{T}
=J^{-1}\mathbf{F}\mathbf{S}\mathbf{F}^{T}.
\label{eq:15}
\end{equation}
The second form shows that $\bm{\sigma}$ is symmetric precisely when $\mathbf{F}^{-1}\mathbf{T}$ is,
so the symmetry condition of the lecture is the referential expression of the familiar symmetry
of the Cauchy stress. For the fluid, eq.~(\ref{eq:12}) gives
\begin{equation}
\bm{\sigma}=J^{-1}\,V'(J)\,J\,\mathbf{F}^{-T}\mathbf{F}^{T}=V'(J)\,\mathbf{I}=-p\,\mathbf{I},\quad p=-V'(J).
\label{eq:16}
\end{equation}
The identification of $p$ as a pressure is thermodynamically sensible: a reference volume
$\dd V$ has energy $V(J)\dd V$ and current volume $\dd v=J\dd V$, so
$p=-\partial(V\dd V)/\partial(\dd v)=-V'(J)$ is minus the derivative of the energy with respect to
volume, and compressing the fluid ($J<1$, $V'<0$) produces a positive pressure. The first
Piola--Kirchhoff stress $\mathbf{T}=-pJ\mathbf{F}^{-T}$ is the same isotropic pressure referred,
through Nanson's formula, to reference rather than current area, and this is why it is not itself
isotropic.

(d) With $\mathbf{F}=\mathbf{I}+s\mathbf{H}$ we have $J=1+s\,\mathrm{tr}\,\mathbf{H}+O(s^{2})$, so
$V'(J)=V'(1)+V''(1)(J-1)+O(s^{2})=s\,\kappa\,\mathrm{tr}\,\mathbf{H}+O(s^{2})$ with $\kappa=V''(1)$,
while $J\mathbf{F}^{-T}=\mathbf{I}+O(s)$. Eq.~(\ref{eq:12}) then gives
\begin{equation}
T_{ij}=s\,\kappa\,\delta_{ij}H_{kk}+O(s^{2})=s\,\kappa\,\delta_{ij}\delta_{kl}H_{kl}+O(s^{2}),
\quad\text{so}\quad A_{ijkl}=\kappa\,\delta_{ij}\delta_{kl},
\label{eq:17}
\end{equation}
which is the isotropic tensor of Example~1 with $\mu=0$, in which case
$\lambda+\frac{2}{3}\mu$ reduces to $\lambda$, and the lecture's $\kappa$ is this remaining modulus.
The name is appropriate: $p=-V'(J)=-\kappa(J-1)+O\big((J-1)^{2}\big)$, so $\kappa=-\ddns p/\ddns J$
at $J=1$ is the bulk modulus in the usual sense. Had the equilibrium not been stress-free, say
$V'(1)=-p^{0}$, the expansion $J\mathbf{F}^{-T}=\mathbf{I}+s(\mathrm{tr}\,\mathbf{H}\,\mathbf{I}-\mathbf{H}^{T})+O(s^{2})$
would have contributed the additional first-order term
$-p^{0}(\delta_{ij}\delta_{kl}-\delta_{il}\delta_{jk})H_{kl}$ to $T_{ij}$, and the resulting
coefficient of $H_{kl}$ does not possess the symmetry $A_{ijkl}=A_{jikl}$. This is a first glimpse of
why the pre-stressed theory needed later in the course is more complicated.
\end{solution}

\begin{example}{The material symmetry group}
Fix a point $\mathbf{x}$ and let $\mathcal{G}$ be the set of rotation matrices $\mathbf{Q}$ such
that $W(\mathbf{F}\mathbf{Q})=W(\mathbf{F})$ for every deformation gradient $\mathbf{F}$.
(a) Show directly from this definition that $\mathcal{G}$ is a group under matrix multiplication.
(b) Show that if $U(\mathbf{C})=f(\mathrm{tr}\,\mathbf{C},\mathrm{tr}\,\mathbf{C}^{2},\det\mathbf{C})$
then every rotation lies in $\mathcal{G}$, and explain why every isotropic strain energy function
is of this form. (c) Let $\hat{\bm{\nu}}$ be a unit vector and
$U(\mathbf{C})=f(\hat{\bm{\nu}}\cdot\mathbf{C}\hat{\bm{\nu}},\hat{\bm{\nu}}\cdot\mathbf{C}^{2}\hat{\bm{\nu}},
\mathrm{tr}\,\mathbf{C},\mathrm{tr}\,\mathbf{C}^{2},\det\mathbf{C})$. Show that
$\mathbf{Q}\in\mathcal{G}$ whenever $\mathbf{Q}\hat{\bm{\nu}}=\pm\hat{\bm{\nu}}$, and conversely
that if $\mathbf{Q}\in\mathcal{G}$ for every $f$ of this form then $\mathbf{Q}\hat{\bm{\nu}}=\pm\hat{\bm{\nu}}$.
(d) Show that if $\mathbf{Q}\in\mathcal{G}$ then the elastic tensor satisfies
$A_{ijkl}=Q_{im}Q_{jn}Q_{kp}Q_{lq}A_{mnpq}$.
\end{example}

\begin{solution}
(a) The identity lies in $\mathcal{G}$ trivially, and associativity is inherited from matrix
multiplication. If $\mathbf{Q}_{1},\mathbf{Q}_{2}\in\mathcal{G}$ then, applying the defining
property twice,
\begin{equation}
W(\mathbf{F}\mathbf{Q}_{1}\mathbf{Q}_{2})=W\big((\mathbf{F}\mathbf{Q}_{1})\mathbf{Q}_{2}\big)
=W(\mathbf{F}\mathbf{Q}_{1})=W(\mathbf{F}),
\label{eq:18}
\end{equation}
so $\mathbf{Q}_{1}\mathbf{Q}_{2}\in\mathcal{G}$; here we used that $\mathbf{F}\mathbf{Q}_{1}$ is
itself an admissible deformation gradient, having positive determinant. If $\mathbf{Q}\in\mathcal{G}$
then $\mathbf{Q}^{-1}=\mathbf{Q}^{T}$ is a rotation and, applying the defining property to the
deformation gradient $\mathbf{F}\mathbf{Q}^{-1}$,
\begin{equation}
W(\mathbf{F}\mathbf{Q}^{-1})=W\big((\mathbf{F}\mathbf{Q}^{-1})\mathbf{Q}\big)=W(\mathbf{F}),
\label{eq:19}
\end{equation}
so $\mathbf{Q}^{-1}\in\mathcal{G}$. Hence $\mathcal{G}$ is a subgroup of the rotation group.

(b) Under $\mathbf{F}\rightarrow\mathbf{F}\mathbf{Q}$ the right Cauchy--Green tensor transforms as
$\mathbf{C}\rightarrow\mathbf{Q}^{T}\mathbf{C}\mathbf{Q}$, and by the cyclic property of the trace
and the multiplicativity of the determinant
\begin{equation}
\mathrm{tr}(\mathbf{Q}^{T}\mathbf{C}\mathbf{Q})=\mathrm{tr}\,\mathbf{C},\quad
\mathrm{tr}\big[(\mathbf{Q}^{T}\mathbf{C}\mathbf{Q})^{2}\big]=\mathrm{tr}(\mathbf{Q}^{T}\mathbf{C}^{2}\mathbf{Q})=\mathrm{tr}\,\mathbf{C}^{2},\quad
\det(\mathbf{Q}^{T}\mathbf{C}\mathbf{Q})=\det\mathbf{C}.
\label{eq:20}
\end{equation}
Thus $U(\mathbf{Q}^{T}\mathbf{C}\mathbf{Q})=U(\mathbf{C})$ for every rotation, which is the
isotropy condition of the lecture. Conversely, the lecture shows that an isotropic $U$ depends on
$\mathbf{C}$ only through its eigenvalues $c_{1},c_{2},c_{3}$. These are the roots of the
characteristic polynomial
\begin{equation}
\begin{split}
&\det(\mathbf{C}-c\mathbf{I})=-c^{3}+I_{1}c^{2}-I_{2}c+I_{3},\\
&I_{1}=\mathrm{tr}\,\mathbf{C},\quad I_{2}=\frac{1}{2}\left[(\mathrm{tr}\,\mathbf{C})^{2}-\mathrm{tr}\,\mathbf{C}^{2}\right],\quad
I_{3}=\det\mathbf{C},
\end{split}
\label{eq:21}
\end{equation}
so the eigenvalues are functions of $(\mathrm{tr}\,\mathbf{C},\mathrm{tr}\,\mathbf{C}^{2},\det\mathbf{C})$
and every isotropic strain energy function can be written in the stated form. The St
Venant--Kirchhoff energy of Example~1 is an instance, with
$\mathrm{tr}\,\mathbf{E}=\frac{1}{2}(I_{1}-3)$ and $\mathrm{tr}(\mathbf{E}^{2})=\frac{1}{4}(\mathrm{tr}\,\mathbf{C}^{2}-2I_{1}+3)$.

(c) Since $(\mathbf{Q}^{T}\mathbf{C}\mathbf{Q})^{2}=\mathbf{Q}^{T}\mathbf{C}^{2}\mathbf{Q}$, the two
new arguments transform as
\begin{equation}
\hat{\bm{\nu}}\cdot(\mathbf{Q}^{T}\mathbf{C}\mathbf{Q})\hat{\bm{\nu}}=(\mathbf{Q}\hat{\bm{\nu}})\cdot\mathbf{C}(\mathbf{Q}\hat{\bm{\nu}}),\quad
\hat{\bm{\nu}}\cdot(\mathbf{Q}^{T}\mathbf{C}\mathbf{Q})^{2}\hat{\bm{\nu}}=(\mathbf{Q}\hat{\bm{\nu}})\cdot\mathbf{C}^{2}(\mathbf{Q}\hat{\bm{\nu}}),
\label{eq:22}
\end{equation}
and both are quadratic in $\mathbf{Q}\hat{\bm{\nu}}$, so they are unchanged when
$\mathbf{Q}\hat{\bm{\nu}}=\pm\hat{\bm{\nu}}$; together with eq.~(\ref{eq:20}) this
shows that such $\mathbf{Q}$ lie in $\mathcal{G}$. For the converse it suffices to consider
$f$ depending only on its first argument, say $U(\mathbf{C})=\hat{\bm{\nu}}\cdot\mathbf{C}\hat{\bm{\nu}}$,
and to choose the particular deformations $\mathbf{C}=\mathbf{I}+\tau\hat{\bm{\nu}}\hat{\bm{\nu}}^{T}$,
which are symmetric and positive-definite for $\tau>-1$. Then
$\hat{\bm{\nu}}\cdot\mathbf{C}\hat{\bm{\nu}}=1+\tau$ while
$(\mathbf{Q}\hat{\bm{\nu}})\cdot\mathbf{C}(\mathbf{Q}\hat{\bm{\nu}})=1+\tau(\hat{\bm{\nu}}\cdot\mathbf{Q}\hat{\bm{\nu}})^{2}$,
and equality for all $\tau$ forces $(\hat{\bm{\nu}}\cdot\mathbf{Q}\hat{\bm{\nu}})^{2}=1$. As
both vectors are unit vectors, the Cauchy--Schwarz inequality then gives
$\mathbf{Q}\hat{\bm{\nu}}=\pm\hat{\bm{\nu}}$. The material symmetry group of a strain energy
function of this form is therefore the set of rotations preserving the axis $\pm\hat{\bm{\nu}}$:
all rotations about $\hat{\bm{\nu}}$, which is the defining property of transverse isotropy in the
lecture, together with the rotations through $\pi$ about axes perpendicular to $\hat{\bm{\nu}}$,
which reverse $\hat{\bm{\nu}}$. That the latter are automatically symmetries reflects the fact
that $\hat{\bm{\nu}}$ enters only quadratically, and correspondingly the transversely isotropic
elastic tensor of the lecture is unchanged by $\hat{\bm{\nu}}\rightarrow-\hat{\bm{\nu}}$.

(d) Material symmetry and material frame indifference together give, for $\mathbf{Q}\in\mathcal{G}$,
\begin{equation}
W(\mathbf{Q}^{T}\mathbf{F}\mathbf{Q})=W(\mathbf{Q}^{T}\mathbf{F})=W(\mathbf{F}),
\label{eq:23}
\end{equation}
where the first equality is material symmetry applied to the deformation gradient
$\mathbf{Q}^{T}\mathbf{F}$ and the second is frame indifference with the rotation $\mathbf{Q}^{T}$.
Writing $G_{mn}=Q_{am}F_{ab}Q_{bn}$ for the components of $\mathbf{Q}^{T}\mathbf{F}\mathbf{Q}$ we have
$\partial G_{mn}/\partial F_{ij}=Q_{im}Q_{jn}$, and differentiating eq.~(\ref{eq:23})
twice by the chain rule,
\begin{equation}
\frac{\partial^{2}W}{\partial F_{ij}\partial F_{kl}}(\mathbf{F})
=Q_{im}Q_{jn}Q_{kp}Q_{lq}\frac{\partial^{2}W}{\partial G_{mn}\partial G_{pq}}(\mathbf{Q}^{T}\mathbf{F}\mathbf{Q}).
\label{eq:24}
\end{equation}
Setting $\mathbf{F}=\mathbf{I}$, for which $\mathbf{Q}^{T}\mathbf{F}\mathbf{Q}=\mathbf{I}$ as well,
gives $A_{ijkl}=Q_{im}Q_{jn}Q_{kp}Q_{lq}A_{mnpq}$. The elastic tensor is thus invariant under the
symmetry group, and this is the form in which material symmetry is used in linearised elasticity;
the point of the argument is that the material symmetry condition alone relates the second
derivatives at $\mathbf{I}$ to those at $\mathbf{Q}$, and frame indifference is needed to bring
them back to $\mathbf{I}$. Numerically, the transversely isotropic tensor of the lecture with
$\hat{\bm{\nu}}=\hat{\mathbf{e}}_{3}$ is invariant under rotations about $\hat{\mathbf{e}}_{3}$ and
under a rotation through $\pi$ about $\hat{\mathbf{e}}_{1}$, but not under a generic rotation
unless $\gamma=\xi=\zeta=0$.
\end{solution}

\begin{example}{Counting components and Voigt notation}
Define the Voigt map from symmetric index pairs to a single index,
\begin{equation}
(11,\,22,\,33,\,23,\,13,\,12)\;\longrightarrow\;(1,\,2,\,3,\,4,\,5,\,6).
\label{eq:25}
\end{equation}
(a) Show that the symmetries $A_{ijkl}=A_{jikl}=A_{ijlk}=A_{klij}$ are exactly what is needed for
the elastic tensor to be represented by a symmetric $6\times6$ matrix $C_{IJ}$, and hence that it
has 21 independent components. (b) Write the isotropic elastic tensor as a $6\times6$ matrix.
(c) For the transversely isotropic tensor of the lecture with $\hat{\bm{\nu}}=\hat{\mathbf{e}}_{3}$,
compute the $6\times6$ matrix and identify the five independent constants
$A_{1111}$, $A_{3333}$, $A_{1133}$, $A_{2323}$ and $A_{1212}$ in terms of $\lambda,\mu,\gamma,\xi,\zeta$.
\end{example}

\begin{solution}
(a) The symmetry $A_{ijkl}=A_{jikl}$ says that $A_{ijkl}$ depends on $(i,j)$ only through the
unordered pair $\{i,j\}$, and $A_{ijkl}=A_{ijlk}$ says the same of $(k,l)$. There are six unordered
pairs, and eq.~(\ref{eq:25}) labels them by $I,J\in\{1,\dots,6\}$, so
\begin{equation}
C_{IJ}=A_{ijkl},\quad I\leftrightarrow\{i,j\},\quad J\leftrightarrow\{k,l\},
\label{eq:26}
\end{equation}
is well defined and every one of the 81 components of $A_{ijkl}$ can be read off from it. The
hyperelastic symmetry $A_{ijkl}=A_{klij}$ is then precisely the statement $C_{IJ}=C_{JI}$. A
symmetric $6\times6$ matrix has $6$ diagonal and $15$ off-diagonal entries, giving the $21$
independent elastic constants quoted in the lecture. For comparison, the hyperelastic symmetry
alone makes the $9\times9$ matrix indexed by ordered pairs $(ij)$ and $(kl)$ symmetric, which gives
the $45$ constants mentioned in the lecture before the remaining symmetries are imposed. When
Hooke's law $T_{ij}=A_{ijkl}e_{kl}$ is written in this notation as $T_{I}=\sum_{J}C_{IJ}\tilde{e}_{J}$
one must take $\tilde{e}_{J}=e_{kl}$ for $J\le3$ but $\tilde{e}_{J}=2e_{kl}$ for $J\ge4$, because both
$(k,l)$ and $(l,k)$ contribute to the sum over the ordered pair; these are the so-called engineering
shear strains.

(b) For $A_{ijkl}=\lambda\delta_{ij}\delta_{kl}+\mu(\delta_{ik}\delta_{jl}+\delta_{il}\delta_{jk})$
the first term contributes $\lambda$ to $C_{IJ}$ whenever $I,J\le3$, and the second contributes
$2\mu$ on the diagonal for $I\le3$ and $\mu$ on the diagonal for $I\ge4$, so
\begin{equation}
\mathbf{C}^{\mathrm{iso}}=\begin{pmatrix}
\lambda+2\mu & \lambda & \lambda & 0 & 0 & 0\\
\lambda & \lambda+2\mu & \lambda & 0 & 0 & 0\\
\lambda & \lambda & \lambda+2\mu & 0 & 0 & 0\\
0 & 0 & 0 & \mu & 0 & 0\\
0 & 0 & 0 & 0 & \mu & 0\\
0 & 0 & 0 & 0 & 0 & \mu
\end{pmatrix}.
\label{eq:27}
\end{equation}

(c) With $\hat{\nu}_{i}=\delta_{i3}$ every factor of $\hat{\nu}$ in the lecture's expression
forces the corresponding index to equal $3$, and the remaining Kronecker deltas must pair equal
indices. Consider first the components with all indices in $\{1,2\}$: only the $\lambda$ and $\mu$
terms survive, giving $A_{1111}=A_{2222}=\lambda+2\mu$, $A_{1122}=\lambda$ and $A_{1212}=\mu$, exactly
as in the isotropic case. For $A_{3333}$ every term contributes: the $\gamma$ term gives $8\gamma$,
the two $\xi$ terms give $4\xi+4\xi$, and each of the four $\zeta$ terms gives $-\zeta$, so
\begin{equation}
A_{3333}=\lambda+2\mu+8\gamma+8\xi-4\zeta.
\label{eq:28}
\end{equation}
For $A_{1133}$ the $\mu$, $\gamma$ and $\zeta$ terms all require a delta between the index $1$ and the
index $3$ and vanish, while the second $\xi$ term $\delta_{ij}\hat{\nu}_{k}\hat{\nu}_{l}$ contributes
$4\xi$, so $A_{1133}=A_{2233}=\lambda+4\xi$. For $A_{2323}$ the $\mu$ term $\delta_{ik}\delta_{jl}$
gives $\mu$, the $\xi$ terms vanish since $\delta_{23}=0$, and of the $\zeta$ terms only
$\hat{\nu}_{j}\hat{\nu}_{l}\delta_{ik}$ survives, giving $A_{2323}=A_{1313}=\mu-\zeta$. Every remaining
entry vanishes: in any non-vanishing term each index is either set equal to $3$ by a factor of
$\hat{\nu}$ or paired with an equal index by a Kronecker delta, so each of the values $1$, $2$ and $3$
must occur an even number of times among $(i,j,k,l)$, and the components computed above, together
with those related to them by the symmetries of part (a), are the only ones with this property.
Collecting the results,
\begin{equation}
\mathbf{C}^{\mathrm{TI}}=\begin{pmatrix}
C_{11} & C_{12} & C_{13} & 0 & 0 & 0\\
C_{12} & C_{11} & C_{13} & 0 & 0 & 0\\
C_{13} & C_{13} & C_{33} & 0 & 0 & 0\\
0 & 0 & 0 & C_{44} & 0 & 0\\
0 & 0 & 0 & 0 & C_{44} & 0\\
0 & 0 & 0 & 0 & 0 & C_{66}
\end{pmatrix},
\label{eq:29}
\end{equation}
with
\begin{equation}
\begin{split}
&C_{11}=A_{1111}=\lambda+2\mu,\quad C_{33}=A_{3333}=\lambda+2\mu+8\gamma+8\xi-4\zeta,\\
&C_{13}=A_{1133}=\lambda+4\xi,\quad C_{44}=A_{2323}=\mu-\zeta,\\
&C_{66}=A_{1212}=\mu,\quad C_{12}=A_{1122}=\lambda=C_{11}-2C_{66}.
\end{split}
\label{eq:30}
\end{equation}
The matrix contains six distinct entries, but $C_{12}$ is fixed by $C_{11}$ and $C_{66}$, so there
are five independent constants, and the map
$(\lambda,\mu,\gamma,\xi,\zeta)\mapsto(C_{11},C_{33},C_{13},C_{44},C_{66})$ is invertible:
$\mu=C_{66}$, $\lambda=C_{11}-2C_{66}$, $\xi=\frac{1}{4}(C_{13}-\lambda)$, $\zeta=\mu-C_{44}$ and
$\gamma=\frac{1}{8}(C_{33}-\lambda-2\mu-8\xi+4\zeta)$. Isotropy is recovered when
$\gamma=\xi=\zeta=0$. For the sample values $\lambda=2$, $\mu=1$, $\gamma=0.3$, $\xi=0.2$ and
$\zeta=0.1$ the script builds all 81 components, confirms the three symmetries of part (a), and
gives $C_{11}=4$, $C_{33}=7.6$, $C_{13}=2.8$, $C_{44}=0.9$, $C_{66}=1$ and $C_{12}=2$, in agreement
with eq.~(\ref{eq:30}). The relation $C_{66}=\frac{1}{2}(C_{11}-C_{12})$ is what makes
the material isotropic within the plane perpendicular to $\hat{\bm{\nu}}$; the five constants are
often written $A=C_{11}$, $C=C_{33}$, $F=C_{13}$, $L=C_{44}$ and $N=C_{66}$ in seismology, and this
is the form in which the anisotropy of the uppermost mantle is parameterised in earth models such as
PREM.
\end{solution}

\begin{example}{From the elastic tensor to the Navier equation}
Insert the isotropic elastic tensor into the linearised equation of motion
$\rho\,\partial^{2}u_{i}/\partial t^{2}=\partial_{j}(A_{ijkl}\partial_{l}u_{k})$. (a) For constant
$\lambda$ and $\mu$ obtain
$\rho\,\ddot{\mathbf{u}}=(\lambda+\mu)\nabla(\nabla\cdot\mathbf{u})+\mu\nabla^{2}\mathbf{u}$.
(b) Show that in a heterogeneous isotropic body additional terms appear, and write them in terms
of $\nabla\lambda$, $\nabla\mu$ and the linearised strain $e_{ij}$.
\end{example}

\begin{solution}
We first form the stress. Contracting the isotropic tensor with the displacement gradient,
\begin{equation}
A_{ijkl}\frac{\partial u_{k}}{\partial x_{l}}
=\lambda\,\delta_{ij}\frac{\partial u_{k}}{\partial x_{k}}+\mu\left(\frac{\partial u_{i}}{\partial x_{j}}+\frac{\partial u_{j}}{\partial x_{i}}\right)
=\lambda\,\delta_{ij}\,e_{kk}+2\mu\,e_{ij},
\label{eq:31}
\end{equation}
which is Hooke's law for an isotropic body, with only the symmetric part of the displacement
gradient entering, as the lecture notes. Taking the divergence, and allowing $\lambda$ and $\mu$
to depend on position,
\begin{equation}
\frac{\partial}{\partial x_{j}}\left(A_{ijkl}\frac{\partial u_{k}}{\partial x_{l}}\right)
=\frac{\partial}{\partial x_{i}}\left(\lambda\frac{\partial u_{k}}{\partial x_{k}}\right)
+\frac{\partial}{\partial x_{j}}\left[\mu\left(\frac{\partial u_{i}}{\partial x_{j}}+\frac{\partial u_{j}}{\partial x_{i}}\right)\right].
\label{eq:32}
\end{equation}

(a) If $\lambda$ and $\mu$ are constant they pass through the derivatives, and using the equality
of mixed partial derivatives in the last term,
\begin{equation}
\frac{\partial}{\partial x_{j}}\left(A_{ijkl}\frac{\partial u_{k}}{\partial x_{l}}\right)
=\lambda\frac{\partial^{2}u_{k}}{\partial x_{i}\partial x_{k}}+\mu\frac{\partial^{2}u_{i}}{\partial x_{j}\partial x_{j}}
+\mu\frac{\partial^{2}u_{k}}{\partial x_{i}\partial x_{k}}
=(\lambda+\mu)\frac{\partial}{\partial x_{i}}(\nabla\cdot\mathbf{u})+\mu\nabla^{2}u_{i},
\label{eq:33}
\end{equation}
so the linearised equation of motion becomes
\begin{equation}
\rho\frac{\partial^{2}\mathbf{u}}{\partial t^{2}}=(\lambda+\mu)\nabla(\nabla\cdot\mathbf{u})+\mu\nabla^{2}\mathbf{u},
\label{eq:34}
\end{equation}
which is the form quoted in the lecture. Using the identity
$\nabla^{2}\mathbf{u}=\nabla(\nabla\cdot\mathbf{u})-\nabla\times(\nabla\times\mathbf{u})$ it can also be
written as $\rho\,\ddot{\mathbf{u}}=(\lambda+2\mu)\nabla(\nabla\cdot\mathbf{u})-\mu\nabla\times(\nabla\times\mathbf{u})$,
which separates the dilatational and rotational parts of the motion; the combinations
$(\lambda+2\mu)/\rho$ and $\mu/\rho$ will reappear as the squared P- and S-wave speeds.

(b) In general the product rule applied to eq.~(\ref{eq:32}) produces, in addition to the
terms of eq.~(\ref{eq:33}) evaluated with the local values of $\lambda$ and $\mu$, the terms
in which the derivative falls on a modulus,
\begin{equation}
\frac{\partial\lambda}{\partial x_{i}}\frac{\partial u_{k}}{\partial x_{k}}
+\frac{\partial\mu}{\partial x_{j}}\left(\frac{\partial u_{i}}{\partial x_{j}}+\frac{\partial u_{j}}{\partial x_{i}}\right)
=(\nabla\cdot\mathbf{u})\frac{\partial\lambda}{\partial x_{i}}+2e_{ij}\frac{\partial\mu}{\partial x_{j}},
\label{eq:35}
\end{equation}
so that the equation of motion in a heterogeneous isotropic body reads
\begin{equation}
\begin{split}
\rho\frac{\partial^{2}u_{i}}{\partial t^{2}}=\,&(\lambda+\mu)\frac{\partial}{\partial x_{i}}(\nabla\cdot\mathbf{u})+\mu\nabla^{2}u_{i}\\
&+(\nabla\cdot\mathbf{u})\frac{\partial\lambda}{\partial x_{i}}+2e_{ij}\frac{\partial\mu}{\partial x_{j}}.
\end{split}
\label{eq:36}
\end{equation}
In vector notation the extra terms are $(\nabla\cdot\mathbf{u})\nabla\lambda+(\nabla\mathbf{u}+\nabla\mathbf{u}^{T})\nabla\mu$.
They are of first order in the gradients of the moduli and first order in the derivatives of
$\mathbf{u}$, whereas the retained terms involve second derivatives of $\mathbf{u}$. For waves whose
wavelength is short compared with the scale over which $\lambda$ and $\mu$ vary they are
therefore small, which is the starting point of ray theory later in the course; but they cannot be
dropped near sharp gradients, and across a discontinuity the divergence form
$\partial_{j}(A_{ijkl}\partial_{l}u_{k})$ must be retained, since it is this form that yields the
continuity of traction as a boundary condition. The script confirms eq.~(\ref{eq:36})
for a quadratic displacement field and linearly varying moduli, for which central differences
of the stress are exact.
\end{solution}

\begin{example}{Uniform dilatation and the bulk modulus}
Consider the uniform dilatation $\mathbf{F}=(1+\epsilon)\mathbf{I}$ of the St Venant--Kirchhoff
material of Example~1. (a) Evaluate $\mathbf{E}$, $\mathbf{S}$ and $\mathbf{T}$ exactly,
and expand $\mathbf{T}$ to second order in $\epsilon$; show that the first-order term is
$\mathbf{T}=3\kappa\epsilon\,\mathbf{I}$ with $\kappa=\lambda+\frac{2}{3}\mu$. (b) Compute the Cauchy
stress $\bm{\sigma}=J^{-1}\mathbf{T}\mathbf{F}^{T}$ and the pressure $p$ to second order in $\epsilon$,
and show that $p=-\kappa(J-1)$ to first order. Give the exact dependence of $p$ on $J$.
\end{example}

\begin{solution}
(a) With $\mathbf{F}=(1+\epsilon)\mathbf{I}$ we have $\mathbf{C}=(1+\epsilon)^{2}\mathbf{I}$ and
$J=(1+\epsilon)^{3}$, so
\begin{equation}
\mathbf{E}=\frac{1}{2}\left[(1+\epsilon)^{2}-1\right]\mathbf{I}=\left(\epsilon+\frac{1}{2}\epsilon^{2}\right)\mathbf{I},\quad
\mathrm{tr}\,\mathbf{E}=3\left(\epsilon+\frac{1}{2}\epsilon^{2}\right).
\label{eq:37}
\end{equation}
From eq.~(\ref{eq:3}) the second Piola--Kirchhoff stress is
$\mathbf{S}=(3\lambda+2\mu)(\epsilon+\frac{1}{2}\epsilon^{2})\mathbf{I}$, and since $3\lambda+2\mu=3\kappa$,
\begin{equation}
\mathbf{T}=\mathbf{F}\mathbf{S}=3\kappa\,(1+\epsilon)\left(\epsilon+\frac{1}{2}\epsilon^{2}\right)\mathbf{I}
=3\kappa\left(\epsilon+\frac{3}{2}\epsilon^{2}+\frac{1}{2}\epsilon^{3}\right)\mathbf{I},
\label{eq:38}
\end{equation}
which is exact. To first order $\mathbf{T}=3\kappa\epsilon\,\mathbf{I}$, in agreement with Hooke's
law for the displacement gradient $H_{kl}=\epsilon\delta_{kl}$, for which the isotropic tensor of
eq.~(\ref{eq:6}) gives $A_{ijkl}H_{kl}=(3\lambda+2\mu)\epsilon\,\delta_{ij}$. The traction on any
reference surface element is $3\kappa\epsilon\,\hat{\mathbf{n}}$ per unit reference area, so an
outward tension is needed to hold the body in its expanded state, as expected.

(b) Since $\mathbf{F}^{T}=(1+\epsilon)\mathbf{I}$ and $J=(1+\epsilon)^{3}$, eq.~(\ref{eq:15}) gives
\begin{equation}
\bm{\sigma}=J^{-1}\mathbf{T}\mathbf{F}^{T}=3\kappa\,\frac{\epsilon+\frac{1}{2}\epsilon^{2}}{1+\epsilon}\,\mathbf{I}
=3\kappa\left(\epsilon-\frac{1}{2}\epsilon^{2}\right)\mathbf{I}+O(\epsilon^{3}),
\label{eq:39}
\end{equation}
and the pressure is
\begin{equation}
p=-\frac{1}{3}\mathrm{tr}\,\bm{\sigma}=-3\kappa\left(\epsilon-\frac{1}{2}\epsilon^{2}\right)+O(\epsilon^{3}).
\label{eq:40}
\end{equation}
At first order $p=-3\kappa\epsilon$, and as $J-1=3\epsilon+O(\epsilon^{2})$ this is $p=-\kappa(J-1)$:
the constant $\kappa$ defined in the lecture through $\lambda$ and $\mu$ is indeed the bulk
modulus, $\kappa=-J\,\ddns p/\ddns J$ evaluated at $J=1$, exactly as for the elastic fluid of
Example~2. The exact result follows on noting that
$(\epsilon+\frac{1}{2}\epsilon^{2})/(1+\epsilon)=\frac{1}{2}\left[(1+\epsilon)-(1+\epsilon)^{-1}\right]$ and
$1+\epsilon=J^{1/3}$, whence
\begin{equation}
p=-\frac{3\kappa}{2}\left(J^{1/3}-J^{-1/3}\right).
\label{eq:41}
\end{equation}
For $\epsilon=0.01$ and $\lambda=2$, $\mu=1.3$ (so that $\kappa=2.8667$) the script gives
$p=-0.085574$, compared with $-3\kappa\epsilon=-0.086$ at first order and
$-3\kappa(\epsilon-\frac{1}{2}\epsilon^{2})=-0.085570$ at second order.

Two remarks. First, the first Piola--Kirchhoff and Cauchy stresses agree at first order in
$\epsilon$, as the lecture asserts for linearised motions about a stress-free equilibrium, but their
second-order terms have coefficients $+\frac{3}{2}$ and $-\frac{1}{2}$ respectively; the difference is
entirely the geometric factor $J^{-1}\mathbf{F}^{T}$ converting force per unit reference area to
force per unit current area, and it is invisible in the linearised theory precisely because the
equilibrium stress vanishes. Had the reference configuration carried a stress $\mathbf{T}^{0}\neq\mathbf{0}$,
this factor would contribute at first order, as in the remark at the end of Example~2. Second, the
strain energy for this deformation is $W=\frac{9}{8}\kappa\,(J^{2/3}-1)^{2}$, which tends to the
finite value $\frac{9}{8}\kappa$ as $J\rightarrow0$, so that crushing the body to zero volume costs a
finite energy. This is a well-known deficiency of the St Venant--Kirchhoff model, which is
intended only for small strains, and it is one reason why the general framework of the lecture is
formulated for an arbitrary function $U(\mathbf{C})$ rather than for this particular choice.
\end{solution}

\end{document}
